وفي سياق التحسنات الجزئية لحساب المعاملات الجارية:
\begin{itemize}
    \item انخفض عجز الميزان التجاري غير البترولي بنحو 390.2 مليون دولار ليبلغ 9.5 مليار دولار، وذلك نتيجة ارتفاع حصيلة الصادرات بمعدل يفوق الزيادة في مدفوعات الواردات.
    \item ارتفع عجز الميزان التجاري البترولي ليصل إلى 5.2 مليار دولار مقابل 4.2 مليار دولار؛ نتيجة زيادة الواردات البترولية.
    \item سجل العجز في ميزان دخل الاستثمار ارتفاعا قدره 2.3\% ليبلغ نحو 4.4 مليار دولار.
\end{itemize}

\vspace{0.4em}
\textbf{المعاملات الرأسمالية والمالية:}

سجل الحساب المالي والرأسمالي صافي تدفق طفيف إلى الخارج في ضوء زيادة الأصول الأجنبية لدى البنوك التجارية، وهو ما عوضه جزئيا تدفقات الاستثمارات الأجنبية المباشرة والاستثمارات في الأوراق المالية:
\begin{itemize}
    \item حقق الاستثمار الأجنبي المباشر صافي تدفق للداخل بلغ نحو 2.4 مليار دولار مقابل 2.7 مليار دولار في الربع السابق.
    \item سجلت الاستثمارات في محفظة الأوراق المالية صافي تدفق للداخل بلغ نحو 1.8 مليار دولار مقابل صافي تدفق للخارج بلغ 384.7 مليون دولار.
    \item سجل التغير في الأصول الأجنبية لدى البنوك صافي تدفق إلى الخارج بلغ نحو 5.3 مليار دولار.
\end{itemize}
